% 来源及取材范围见分题文件；此为整理者转录的正文，不是下载原件。
\begin{definition}[10.84.1]
Let $M$ be an $R$-module. A direct sum d\'evissage of $M$ is a family of submodules $(M_\alpha)_{\alpha\in S}$, indexed by an ordinal $S$ and increasing (with respect to inclusion), such that:
\begin{enumerate}[label=(\arabic*),start=0]
\item $M_0=0$;
\item $M=\bigcup_\alpha M_\alpha$;
\item if $\alpha\in S$ is a limit ordinal, then $M_\alpha=\bigcup_{\beta<\alpha}M_\beta$;
\item if $\alpha+1\in S$, then $M_\alpha$ is a direct summand of $M_{\alpha+1}$.
\end{enumerate}
If moreover
\begin{enumerate}[label=(\arabic*),start=4]
\item $M_{\alpha+1}/M_\alpha$ is countably generated for $\alpha+1\in S$,
\end{enumerate}
then $(M_\alpha)_{\alpha\in S}$ is called a Kaplansky d\'evissage of $M$.
\end{definition}
\begin{lemma}[10.84.2]
Let $M$ be an $R$-module. If $(M_\alpha)_{\alpha\in S}$ is a direct sum d\'evissage of $M$, then $M\cong\bigoplus_{\alpha+1\in S}M_{\alpha+1}/M_\alpha$.
\end{lemma}
\begin{proof}
By property (3) of a direct sum d\'evissage, there is an inclusion $M_{\alpha+1}/M_\alpha\to M$ for each $\alpha\in S$. Consider the map
\[f:\bigoplus_{\alpha+1\in S}M_{\alpha+1}/M_\alpha\to M\]
given by the sum of these inclusions. Further consider the restrictions
\[f_\beta:\bigoplus_{\alpha+1\leq\beta}M_{\alpha+1}/M_\alpha\longrightarrow M\]
for $\beta\in S$. Transfinite induction on $S$ shows that the image of $f_\beta$ is $M_\beta$. For $\beta=0$ this is true by (0). If $\beta+1$ is a successor ordinal and it is true for $\beta$, then it is true for $\beta+1$ by (3). And if $\beta$ is a limit ordinal and it is true for $\alpha<\beta$, then it is true for $\beta$ by (2). Hence $f$ is surjective by (1).

Transfinite induction on $S$ also shows that the restrictions $f_\beta$ are injective. For $\beta=0$ it is true. If $\beta+1$ is a successor ordinal and $f_\beta$ is injective, then let $x$ be in the kernel and write $x=(x_{\alpha+1})_{\alpha+1\leq\beta+1}$ in terms of its components $x_{\alpha+1}\in M_{\alpha+1}/M_\alpha$. By property (3) and the fact that the image of $f_\beta$ is $M_\beta$ both $(x_{\alpha+1})_{\alpha+1\leq\beta}$ and $x_{\beta+1}$ map to $0$. Hence $x_{\beta+1}=0$ and, by the assumption that the restriction $f_\beta$ is injective also $x_{\alpha+1}=0$ for every $\alpha+1\leq\beta$. So $x=0$ and $f_{\beta+1}$ is injective. If $\beta$ is a limit ordinal consider an element $x$ of the kernel. Then $x$ is already contained in the domain of $f_\alpha$ for some $\alpha<\beta$. Thus $x=0$ which finishes the induction. We conclude that $f$ is injective since $f_\beta$ is for each $\beta\in S$.
\end{proof}
\begin{lemma}[10.84.3]
Let $M$ be an $R$-module. Then $M$ is a direct sum of countably generated $R$-modules if and only if it admits a Kaplansky d\'evissage.
\end{lemma}
\begin{proof}
The lemma takes care of the ``if'' direction. Conversely, suppose $M=\bigoplus_{i\in I}N_i$ where each $N_i$ is a countably generated $R$-module. Well-order $I$ so that we can think of it as an ordinal. Then setting $M_i=\bigoplus_{j<i}N_j$ gives a Kaplansky d\'evissage $(M_i)_{i\in I}$ of $M$.
\end{proof}
\begin{theorem}[10.84.4]
Suppose $M$ is a direct sum of countably generated $R$-modules. If $P$ is a direct summand of $M$, then $P$ is also a direct sum of countably generated $R$-modules.
\end{theorem}
\begin{proof}
Write $M=P\oplus Q$. We are going to construct a Kaplansky d\'evissage $(M_\alpha)_{\alpha\in S}$ of $M$ which, in addition to the defining properties (0)--(4), satisfies:
\begin{enumerate}[label=(\arabic*),start=5]
\item Each $M_\alpha$ is a direct summand of $M$;
\item $M_\alpha=P_\alpha\oplus Q_\alpha$, where $P_\alpha=P\cap M_\alpha$ and $Q_\alpha=Q\cap M_\alpha$.
\end{enumerate}
(Note: if properties (0)--(2) hold, then in fact property (3) is equivalent to property (5).)

To see how this implies the theorem, it is enough to show that $(P_\alpha)_{\alpha\in S}$ forms a Kaplansky d\'evissage of $P$. Properties (0), (1), and (2) are clear. By (5) and (6) for $(M_\alpha)$, each $P_\alpha$ is a direct summand of $M$. Since $P_\alpha\subset P_{\alpha+1}$, this implies $P_\alpha$ is a direct summand of $P_{\alpha+1}$; hence (3) holds for $(P_\alpha)$. For (4), note that
\[M_{\alpha+1}/M_\alpha\cong P_{\alpha+1}/P_\alpha\oplus Q_{\alpha+1}/Q_\alpha,\]
so $P_{\alpha+1}/P_\alpha$ is countably generated because this is true of $M_{\alpha+1}/M_\alpha$.

It remains to construct the $M_\alpha$. Write $M=\bigoplus_{i\in I}N_i$ where each $N_i$ is a countably generated $R$-module. Choose a well-ordering of $I$. By transfinite recursion we are going to define an increasing family of submodules $M_\alpha$ of $M$, one for each ordinal $\alpha$, such that $M_\alpha$ is a direct sum of some subset of the $N_i$.

For $\alpha=0$ let $M_0=0$. If $\alpha$ is a limit ordinal and $M_\beta$ has been defined for all $\beta<\alpha$, then define $M_\alpha=\bigcup_{\beta<\alpha}M_\beta$. Since each $M_\beta$ for $\beta<\alpha$ is a direct sum of a subset of the $N_i$, the same will be true of $M_\alpha$. If $\alpha+1$ is a successor ordinal and $M_\alpha$ has been defined, then define $M_{\alpha+1}$ as follows. If $M_\alpha=M$, then let $M_{\alpha+1}=M$.

If not, choose the smallest $j\in I$ such that $N_j$ is not contained in $M_\alpha$. We will construct an infinite matrix $(x_{mn})$, $m,n=1,2,3,\ldots$ such that:
\begin{enumerate}
\item $N_j$ is contained in the submodule of $M$ generated by the entries $x_{mn}$;
\item if we write any entry $x_{k\ell}$ in terms of its $P$- and $Q$-components, $x_{k\ell}=y_{k\ell}+z_{k\ell}$, then the matrix $(x_{mn})$ contains a set of generators for each $N_i$ for which $y_{k\ell}$ or $z_{k\ell}$ has nonzero component.
\end{enumerate}
Then we define $M_{\alpha+1}$ to be the submodule of $M$ generated by $M_\alpha$ and all $x_{mn}$; by property (2) of the matrix $(x_{mn})$, $M_{\alpha+1}$ will be a direct sum of some subset of the $N_i$. To construct the matrix $(x_{mn})$, let $x_{11},x_{12},x_{13},\ldots$ be a countable set of generators for $N_j$. Then if $x_{11}=y_{11}+z_{11}$ is the decomposition into $P$- and $Q$-components, let $x_{21},x_{22},x_{23},\ldots$ be a countable set of generators for the sum of the $N_i$ for which $y_{11}$ or $z_{11}$ have nonzero component. Repeat this process on $x_{12}$ to get elements $x_{31},x_{32},\ldots$, the third row of our matrix. Repeat on $x_{21}$ to get the fourth row, on $x_{13}$ to get the fifth, and so on, going down along successive anti-diagonals as indicated below:
\[\left(\vcenter{\xymatrix@R=2mm@C=2mm{x_{11}&x_{12}\ar[dl]&x_{13}\ar[dl]&x_{14}\ar[dl]&\ldots\\x_{21}&x_{22}\ar[dl]&x_{23}\ar[dl]&\ldots\\x_{31}&x_{32}\ar[dl]&\ldots\\x_{41}&\ldots\\\ldots}}\right).\]
Transfinite induction on $I$ (using the fact that we constructed $M_{\alpha+1}$ to contain $N_j$ for the smallest $j$ such that $N_j$ is not contained in $M_\alpha$) shows that for each $i\in I$, $N_i$ is contained in some $M_\alpha$. Thus, there is some large enough ordinal $S$ satisfying: for each $i\in I$ there is $\alpha\in S$ such that $N_i$ is contained in $M_\alpha$. This means $(M_\alpha)_{\alpha\in S}$ satisfies property (1) of a Kaplansky d\'evissage of $M$.

The family $(M_\alpha)_{\alpha\in S}$ moreover satisfies the other defining properties, and also (5) and (6) above: properties (0), (2), (4), and (6) are clear by construction; property (5) is true because each $M_\alpha$ is by construction a direct sum of some $N_i$; and (3) is implied by (5) and the fact that $M_\alpha\subset M_{\alpha+1}$.
\end{proof}
\begin{theorem}[10.84.5]
If $P$ is a projective $R$-module, then $P$ is a direct sum of countably generated projective $R$-modules.
\end{theorem}
\begin{proof}
A module is projective if and only if it is a direct summand of a free module, so this follows from Theorem 10.84.4.
\end{proof}
